\section{Strong-Comparator and Open3DSG Sensitivities}

\paragraph{Matched strong comparators.}
To separate score construction and model capacity, rank-average, RRF,
and a small two-hidden-unit MLP use the same family-aware ranking procedure,
support/contact source order, public/full 548-context target, and scan-cluster
resampling as the \method{} product. The MLP uses the same training rows,
constructed target, feature contract, exclusion of the source relation score, and split as the
linear heads; it is fitted once and applied to all three predictors without
target-specific refitting.

Let $q_i^Z$ and $q_i^C$ denote descending within-context percentile ranks of
the source relation score and compatibility, and let $r_i^Z,r_i^C$ be their
one-indexed ranks. Rank-average and fixed-constant Reciprocal Rank Fusion use
\begin{equation}
S_i^{\mathrm{rank}}=\tfrac12(q_i^Z+q_i^C),\qquad
S_i^{\mathrm{RRF}}=\frac{1}{60+r_i^Z}+\frac{1}{60+r_i^C}.
\end{equation}

\begin{table}[H]
\centering
\scriptsize
\setlength{\tabcolsep}{3pt}
\caption{\textbf{Matched fusion comparisons.} Cells are Recall / verifier V
at $K=50$ and $K=100$. Every row preserves the source family sequence and
support/contact selections.}
\label{tab:nonlinear-comparison}
\begin{tabular}{llrr}
\toprule
Source & Fusion & $K=50$ & $K=100$ \\
\midrule
\multirow{4}{*}{VL-SAT}
 & Product & .9277/.0197 & .9658/.0295 \\
 & Matched MLP & .9272/.0189 & .9650/.0296 \\
 & Rank-average & .9026/.0162 & .9700/.0225 \\
 & RRF & .9174/.0180 & .9710/.0247 \\
\midrule
\multirow{4}{*}{Open3DSG}
 & Product & .4418/.0342 & .5692/.0324 \\
 & Matched MLP & .4670/.0413 & .5989/.0371 \\
 & Rank-average & .4323/.0458 & .5408/.0548 \\
 & RRF & .4237/.0948 & .5476/.0772 \\
\midrule
\multirow{4}{*}{SGFN}
 & Product & .7450/.0263 & .9303/.0350 \\
 & Matched MLP & .7457/.0258 & .9288/.0350 \\
 & Rank-average & .7022/.0236 & .9280/.0259 \\
 & RRF & .6813/.0253 & .8774/.0302 \\
\bottomrule
\end{tabular}
\end{table}

At $K=50$, the matched MLP has higher Open3DSG Recall than the product but also
higher V; product and MLP are nearly tied on VL-SAT and SGFN. Rank fusion can
lower V at large budgets but incurs larger predictor-dependent Recall losses.
No fusion dominates the joint trade-off. A separate exact-label MLP uses
stronger SGFN-specific correctness supervision and is therefore retained only
as a source-specific upper comparator, not as an apples-to-apples baseline.

\paragraph{Family-conditioning ablation.}
The pooled product fits one compatibility model across all families and does
not use the primary family-aware ranking procedure. Compared with Product (all
families), it can increase Recall but generally raises V,
particularly for Open3DSG and SGFN:

\begin{table}[H]
\centering
\scriptsize
\setlength{\tabcolsep}{2.4pt}
\caption{\textbf{Family-specific versus pooled all-family products.} Cells
are Recall / verifier V at $K=50$ and $K=100$.}
\label{tab:pooled-ablation}
\begin{tabular}{lrrrr}
\toprule
& \multicolumn{2}{c}{$K=50$} & \multicolumn{2}{c}{$K=100$} \\
\cmidrule(lr){2-3}\cmidrule(lr){4-5}
Source & Family & Pooled & Family & Pooled \\
\midrule
VL-SAT & .9293/.0203 & .9300/.0219 & .9688/.0325 & .9690/.0387 \\
Open3DSG & .4655/.0265 & .4675/.0510 & .6005/.0330 & .6402/.0743 \\
SGFN & .7706/.0256 & .7925/.0292 & .9418/.0372 & .9413/.0464 \\
\bottomrule
\end{tabular}
\end{table}

\paragraph{Open3DSG coverage sensitivity.}
The public preprocessing drops a context when fewer than four annotated
objects retain view metadata. This accounts for all 15 missing contexts.
Table~\ref{tab:open3dsg-routes} distinguishes the eligible-context
evaluation, the conservative full-target main evaluation, and a documented
full-coverage recovery.

\begin{table}[H]
\centering
\scriptsize
\setlength{\tabcolsep}{3pt}
\caption{\textbf{Open3DSG coverage sensitivity at $K=100$.} Cells are Recall /
verifier V. The main evaluation uses the full 3,972-relation denominator and no
predictions for the 15 public-pipeline drops.}
\label{tab:open3dsg-routes}
\begin{tabular}{lrrr}
\toprule
Evaluation & Contexts / GT & Source & \method{} \\
\midrule
Public eligible & 533 / 3,899 & .5206/.1242 & .5799/.0324 \\
Public, full target (main) & 548 / 3,972 & .5111/.1242 & .5692/.0324 \\
Recovered full coverage & 548 / 3,972 & .5161/.1242 & .5743/.0332 \\
\bottomrule
\end{tabular}
\end{table}

The conclusion is stable to all three denominator and coverage choices. The
main paper uses the conservative full-target public evaluation; recovered results
are sensitivity evidence rather than the headline benchmark.
